\subsection{The integral sign removes the primitive-factor exception}
\label{sec:primitive-integral-sign}

\paragraph{Finding A46 (the rational primitive criterion in all characteristics).}
Let $C/k$ be a smooth projective geometrically connected curve of
genus $g\ge2$ over a finite field of characteristic $p$.
Fix $P\in C(k)$. Suppose that, over a finite extension $L/k$,
there is a geometrically connected finite etale cover of degree $n$
\[
               D\longrightarrow C_L
\]
which is geometrically Galois with cyclic deck group, where
$n>1$ is odd and $p\nmid n$, and a point $P'\in D(L)$ above $P$.
The geometric deck group need not be constant over $L$.
Let $P_n\subset J_D$ be the actual primitive conductor-$n$
abelian subvariety, defined over $L$ by the integral projector
specified below. If
\[
                         P_n(L)[p^\infty]=0,
 \tag{IS.1}
\]
then the point section of $P$ is distinct from that of every
different point of $C(k)$ in the geometrically pro-prime-to-$p$
arithmetic fundamental group. In fact the same criterion over
$L$ separates $P$ from every distinct point of $C(L)$.

Consequently, if a geometrically cyclic etale cover of odd degree
prime to $p$ has primitive factor of geometric $p$-rank zero,
then it gives the point-section separation conclusion for the
base curve after the usual descent and pointing of the cover.
This geometric sufficient condition applies in every genus
$g\ge2$ and every positive characteristic.
No hypothesis about an elliptic quotient or ordinarity is needed
in either assertion.

\paragraph{Prior result and exact strengthening.}
The recovered \texttt{ANABELIAN\_SUCCESSOR} final row A.14,
working identifier AC.11, supplies the composite-projector
argument with an exclusion of a separable double map of the base
to a genus-one curve. Findings A15 and A18 retained the complete
small-field argument, and removed that exclusion in specified
geometric cases. The new step here uses the signs of the
\emph{integral principal divisor} to show that the degree-two
outcome itself is impossible. It removes the rational-only
qualification in A15/A18 and does not use their
Deuring--Shafarevich computation. We give the full proof, including
the descent of the actual primitive subvariety, independently of
any nonconstant-deck statement elsewhere in this report.

\begin{proof}
\emph{1. The canonical integral primitive projector.}
Over an algebraic closure of $L$, choose a generator $\sigma$
of the cyclic deck group $G$. In $\mathbf Q[G]$, let $e_n$
be the idempotent which is one on the faithful characters and
zero on the characters of proper conductor. Explicitly,
\[
 e_n=\frac1n\sum_{a=0}^{n-1}c_n(a)\sigma^a,
 \qquad
 c_n(a)=\sum_{d\mid\gcd(n,a)}d\,\mu(n/d)\in\mathbf Z.
 \tag{IS.2}
\]
Here the formula for $a=0$ uses $\gcd(n,0)=n$.
To verify the formula, write $\zeta$ for a primitive complex
$n$th root. Inclusion--exclusion on the units modulo $n$ gives
\[
 \sum_{u\in(\mathbf Z/n)^\times}\zeta^{-au}
       =\sum_{d\mid\gcd(n,a)}d\,\mu(n/d)=c_n(a).
\]
The finite Fourier sum therefore makes $e_n$ the stated
idempotent. Put
\[
        p_n=ne_n\in\mathbf Z[G],\qquad
        P_n=\operatorname{im}(p_n:J_D\longrightarrow J_D).
 \tag{IS.3}
\]
The coefficients satisfy $c_n(ua)=c_n(a)$ for every unit $u$
modulo $n$. Arithmetic Galois acts on $G$ by such automorphisms,
so $p_n$ is an endomorphism defined over $L$ and its image is an
abelian subvariety defined over $L$. This is an actual image,
not an unspecified isogenous representative.

Choose the geometric group-ring operator
\[
                 A=\prod_{\ell\mid n}(1-\sigma^{n/\ell}),
 \tag{IS.4}
\]
where the product is over distinct prime divisors. On the
character $\sigma\mapsto\zeta^i$, it acts by
$\prod_{\ell\mid n}(1-\zeta^{i n/\ell})$, nonzero exactly
when $\gcd(i,n)=1$. Hence
\[
                  Ae_n=A,\qquad Ap_n=nA
                       \quad\hbox{in }\mathbf Z[G].
 \tag{IS.5}
\]
The rational identity is integral because $\mathbf Z[G]$
is torsion-free. The images of $A$ and $p_n$ over the algebraic
closure are the same abelian subvariety: $A$ vanishes on every
proper conductor and is invertible on the faithful factor of
$\mathbf Q[G]$. Thus (IS.3) agrees with the primitive factor
used in the constant-deck arguments.

\emph{2. A collision gives a principal reduced projector divisor.}
Suppose that $Q\ne P$ had the same point section, and restrict
the collision to $L$. The covering formulation gives a point
$Q'\in D(L)$ above $Q$, with the same lifted section as $P'$,
and therefore
\[
                        \delta=[Q'-P']
                              \in J_D(L)[p^\infty].
 \tag{IS.6}
\]
For completeness, a finite etale cover which is geometrically
cyclic of degree prime to $p$ is a permitted neighbourhood of
the geometrically pro-prime-to-$p$ group. Equality of sections
lifts to the same pointed neighbourhood. The subgroup topology
is compatible: taking the core of a geometric prime-to-$p$
quotient of the normal index-$n$ subgroup gives a quotient
whose order divides a product of prime-to-$p$ orders and $n$.
The lifted sections have equal prime-to-$p$ abelian Kummer
classes. On a finite-field Jacobian the kernel of that Kummer
map on rational points is exactly the $p$-primary subgroup,
which gives (IS.6). This is the same covering/Kummer implication
used in A1 and the recovered B.13--B.14 arguments.

Now $p_n\delta$ belongs to $P_n(L)[p^\infty]$, so (IS.1)
makes it zero. Equivalently, multiplication by $n$ is an
automorphism on the finite $p$-primary group and
\[
 b=e_n\delta:=[n]^{-1}(p_n\delta)
                      \in P_n(L)[p^\infty]
 \tag{IS.7}
\]
is well-defined and zero. Equation (IS.5) gives
$nA\delta=Ap_n\delta=0$. Since $A\delta$ is $p$-primary and
$p\nmid n$, it follows that
\[
                               A\delta=0.
 \tag{IS.8}
\]
We now use (IS.8) over the algebraic closure. Although $A$ need
not descend to $L$, no constant-field extension has been used
to preserve or replace the rational condition (IS.1).

Put $s=\omega(n)$ and $d=2^s$. The $2^s$ subset products of the
elements $\sigma^{n/\ell}$ are distinct: their factors have
orders in distinct prime-primary components of the cyclic
group. Thus the expansion of $A$ has exactly $d$ coefficients,
each $1$ or $-1$. Its augmentation is zero, so half have each
sign. The two free deck orbits of $P'$ and $Q'$ are disjoint,
because their images on $C$ are distinct. Consequently
$A(Q'-P')$ is a nonzero integral divisor with exactly $d$
simple zeros and $d$ simple poles. Equation (IS.8) makes it
principal. Over $\overline L$ choose $f$ with
\[
                       \operatorname{div}(f)=A(Q'-P').
 \tag{IS.9}
\]
It has degree $d$ and is separable in every characteristic:
a $p$th power has all divisor multiplicities divisible by $p$,
whereas (IS.9) has simple zeros and poles.

\emph{3. The field generated by the deck conjugates.}
Write
\[
 K=\overline L(D),\quad
 M=\overline L(\gamma f:\gamma\in G),\quad e=[K:M].
\]
Every conjugate has degree $d$. The Castelnuovo--Severi
consequence in Poonen, Proposition~2.1, Lemma~2.2 and
Corollary~2.3, gives
\[
                    e\mid d,
                    \qquad g(M)\le(d/e-1)^2.
 \tag{IS.10}
\]
The exact applicability is Corollary~2.3: over the perfect field
$\overline L$, the subfield $M$ is generated by finitely many
elements all of degree $d$, so it is $d$-controlled in the
definition preceding Lemma~2.2. Thus there is a positive integer
$u$ with $[K:M]=d/u$ and $g(M)\le(u-1)^2$; substituting
$u=d/e$ gives (IS.10)
\cite[Proposition~2.1, Lemma~2.2 and Corollary~2.3]{Poonen2007Gonality}.

For every odd $n>1$,
\[
                         n>(2^{\omega(n)}-1)^2=(d-1)^2.
 \tag{IS.11}
\]
For $s=1,2$ use $n\ge3,15$, respectively. For $s\ge3$,
$n\ge3\cdot5^{s-1}>4^s>(2^s-1)^2$; the middle inequality
holds at $s=3$ and its ratio increases with $s$.
Since the original cover is etale,
$g(D)=n(g(C)-1)+1\ge n+1$.
Thus (IS.10)--(IS.11) exclude $e=1$, and $e>1$ is a power of two.

Let $Z$ be the smooth projective curve with function field $M$,
and let $v:D_{\overline L}\to Z$ be the corresponding morphism
of degree $e$. The group $G$ preserves $M$ and acts faithfully
on $Z$: a kernel would have order dividing both $n$ and $e$.
This action is free as well. Indeed, a nontrivial stabilizer
contains an element of some odd prime order $\ell\mid n$.
That element acts freely on $D$ and preserves the fibre divisor
$v^*z$ above the point it fixes on $Z$. Its orbits in that fibre
have length $\ell$, with constant multiplicities along each
orbit, forcing $\ell\mid e$. This contradicts $e$ being a
power of two.

\emph{4. A reduced fibre and its integral signs give the contradiction.}
Write $f=h\circ v$ with $h\in\overline L(Z)$.
Choose a zero $z$ of $h$. The fibre divisor $v^*z$ is reduced:
every multiplicity in it, multiplied by the positive order of
$h$ at $z$, is a zero multiplicity of $f$, which is one.
It therefore consists of $e$ distinct geometric points, all
lying in the two original deck orbits of $P',Q'$.
Each such orbit maps injectively to $Z$, because the induced
$G$-action there is free. Hence $e\le2$, and so $e=2$.

The extension $K/M$ is separable, since it is intermediate in
the separable extension $K/\overline L(f)$. Its quadratic
involution $\tau$ is normalized by $G$, which preserves $M$.
The group of order two has no nontrivial automorphism, so
$\tau$ commutes with $G$.
The reduced fibre just chosen contains one point of each
original orbit. Its two points are exchanged by $\tau$;
therefore $\tau$ exchanges the two entire $G$-orbits.
For some $j$ modulo $n$,
\[
                  \tau P'=\sigma^jQ',\qquad
                  \tau Q'=\sigma^{-j}P'.
 \tag{IS.12}
\]
The second equality follows from the first, $\tau^2=1$,
and commutation with $G$.

Since $f\in M$, its divisor is $\tau$-invariant. Applying
$\tau$ to (IS.9), and comparing its coefficients on the free
$P'$-orbit, gives
\[
                         A\sigma^{-j}=-A
                                \quad\hbox{in }\mathbf Z[G].
 \tag{IS.13}
\]
This comparison is literal: the map from $\mathbf Z[G]$
to the integral divisors supported on that orbit,
$\sum a_\gamma\gamma\mapsto\sum a_\gamma\gamma P'$,
is injective. If $t$ is the order of $\sigma^j$, then $t$
is odd. Iterating (IS.13) gives
\[
                           A=(-1)^tA=-A.
\]
But $\mathbf Z[G]$ is torsion-free and $A\ne0$, as its
subset-product expansion already shows. This is impossible.
The argument uses integer divisor coefficients, so the two
signs remain distinct also in characteristic two.

The assumed collision is therefore excluded. In particular,
there is no need to calculate the genus of $Z$, to descend an
elliptic quotient, or to impose an additional condition on it.

\emph{5. The geometric consequence and its constant-field scope.}
A primitive factor of geometric $p$-rank zero has no nonzero
geometric points of $p$-power order. Thus its rational
$p$-primary subgroup vanishes over every finite constant
extension. Descend the geometric cyclic cover and its deck
action to a finite field, and extend constants once more to
rationalize a chosen point above the prescribed base point.
Condition (IS.1) then holds, so the result just proved applies.
For an initially rational criterion, in contrast, the proof
uses its specified field $L$ throughout the arithmetic step;
it does not enlarge that field to make the deck group constant.
\end{proof}

\paragraph{Effect on the earlier gates.}
A15's rational primitive-factor criterion no longer needs
rank one or exclusion of an elliptic double quotient. A18's
geometric criterion extends to every characteristic for odd
cyclic degrees prime to that characteristic. In characteristic
two and genus two, the already proved actual-theta orbit
coverage test still supplies geometric primitive rank zero,
and hence remains a sufficient input. The proof above does
not assert existence of such a cover or such theta incidence
for every prescribed curve. The closure here is the earlier
\emph{rational-only primitive-factor qualification}, not the
universal incidence gate.
